% ============================================================
% 分片：s03b-genus1  对应原文 PDF 第 27--42 页（印刷页 27--42）
% 负责章节：3.6--3.11（§3 续片，本身不含 \section 标题）
% 译者：trans-s03b   审校状态：初稿
% ============================================================

% ==== 原文 PDF 第 27 页（印刷页 27）====

\subsection{Langlands 对偶群}

$G=SL_2$ 的 Langlands 对偶群为 ${}^{L}G=SO_3$，等价地也为 $PGL_2$。因此我们也希望理解 $C$ 上结构群为 $SO_3$ 的 Higgs 丛。它们与 $SL_2$ Higgs 丛密切相关，因为 $SO_3=SL_2/\mathbb{Z}_2$ 是 $SL_2$ 的伴随形式。

我们用 $W$ 表示一个全纯 $SO_3$ 丛，即带有非退化全纯二次型与体积形式的秩三全纯丛。由于 $SO_3$ 的三维表示就是伴随表示，我们无须区分 $W$ 与相应的伴随丛。

$SO_3$ Higgs 丛 $(W,\varphi)$ 在 Riemann 面 $C$ 上的模空间有两个分支，它们由底丛 $W$ 的第二 Stiefel--Whitney 类 $w_2(W)$ 区分。\footnote{更精确的陈述见注 3.3。} 我们把这些分支记作 $\mathcal{M}_H(C;SO_3,w_2)$，其中 $w_2$ 或为 $0$，或为 $H^2(C,\mathbb{Z}_2)\cong\mathbb{Z}_2$ 的非零元，该非零元我们记作 $\theta$。先考虑 $w_2(W)=0$ 的情形。
% [原文疑误: p.27 原文将 "Stiefel--Whitney" 拼作 "Stieffel-Whitney"（Stiefel 只有一个 f）；译文按标准拼写处理，特此说明。]

$w_2=0$ 的 $SO_3$ 丛 $W$ 的结构群可以提升到 $SL_2$，于是可构造一个带 $SL_2$ 结构群的相伴秩二丛 $E$。$E$ 在相差一个 2 阶线丛的扭转下唯一确定。因此，$w_2$ 为零的 $SO(3)$ Higgs 丛的模空间 $\mathcal{M}_H(C;SO_3,0)$ 与相应的 $SL_2$ 模空间 $\mathcal{M}_H(C;SL_2)$ 之间有非常简单的关系：
\begin{equation}\tag{3.20}
\mathcal{M}_H(C;SO_3,0)=\mathcal{M}_H(C;SL_2)/Q.
\end{equation}

利用我们对 $Q$ 作用的知识，可以立即描述 $\mathcal{M}_H(C;SO_3,0)$ 的 Hitchin 纤维化在特殊点 $w=e_1,e_2,e_3$ 处的纤维。例如在 $w=e_1$ 处，除以 $T_2$（或 $T_3$）的效果是把两个分支 $\boldsymbol{F}^\pm_{e_1}$ 等同起来。因此我们只需关注其中之一，比如 $\boldsymbol{F}^+_{e_1}$。它是一条亏格 $0$ 的曲线，其中两个点被等同（即被 $T_2$ 交换的那两个二重点）。我们还须再除以 $T_1$；商仍是一条带一个二重点的亏格 $0$ 曲线。这就是 ${}^{L}G=SO_3$ 的 Hitchin 纤维化的例外纤维的本性。它们显示在原书 p.18 的右图中。特别地，$SO_3$ 的特殊 Hitchin 纤维不可约，而 $SL_2$ 的相应纤维有两个分支。

现在来看 $SL_2$ 与 $SO_3$ 之间另一个关键区别。在 $SL_2$ 的情形，尽管 Hitchin 纤维化的某些纤维是奇异的，纤维的奇点并不是 $\mathcal{M}_H$ 的奇点；$\mathcal{M}_H$ 在例外纤维附近光滑（事实上，对一般的 $\sigma$，处处光滑）。

% ==== 原文 PDF 第 28 页（印刷页 28）====

然而用 $Q=\mathbb{Z}_2\times\mathbb{Z}_2$ 对 $\mathcal{M}_H(C;SL_2)$ 取商会产生奇点。$T_2$ 在纤维 $\boldsymbol{F}_{e_1}$ 上自由作用，因而在其小邻域上也自由作用。但 $T_1$ 固定纤维的两个二重点。因此除以 $T_1$ 产生两个 $A_1$ 奇点。它们被 $T_2$ 的作用交换，于是商 $\mathcal{M}_H(C;SO_3,0)$ 在每条例外纤维上有一个 $A_1$ 奇点。

\subsubsection{与几何内窥的关系}

我们刚才描述的正是 §2.2 所预告的图景。对 $SO_3$，Hitchin 纤维化的特殊纤维在一点 $r$ 处有一个 $A_1$ 奇点。支撑在一般点 $r^*$ 处的 $B$-膜 $\mathcal{B}^*$ 不可约，但当 $r^*=r$ 时它可以分裂为和 $\mathcal{B}=\mathcal{B}_+\oplus\mathcal{B}_-$。对偶地，$SL_2$ 的 Hitchin 纤维化的特殊纤维有两个分支，因此支撑在该纤维上的 $A$-膜可以分裂为两个 $A$-膜之和，各支撑在一个分支上。由于这些分支单连通，支撑在其中一个分支上的秩 $1$ $A$-膜没有模。这与分数膜 $\mathcal{B}_+$ 与 $\mathcal{B}_-$ 没有模这一事实对偶。但膜 $\mathcal{B}=\mathcal{B}_+\oplus\mathcal{B}_-$ 有模（因为它可以离开奇点作形变），而对偶地，两个 $A$-膜之和 $\mathcal{A}$ 同样可以形变。
% [Tutor Note: 点拨镜像字典的方向。$\mathcal{B}^*$ 是支撑在 $\mathcal{M}_H(SO_3)$ 一点上的点层，在逐纤维 T-对偶下对应于对偶 Hitchin 纤维上的一个秩一平坦酉线丛；$r^*$ 正是该线丛的模（其和乐由 $r^*$ 在对偶环面中的位置给出）。$\mathcal{B}_+$ 与 $\mathcal{B}_-$ 分别对应轨形点处 $\mathbb{Z}_2$ 的两个一维表示，而直和 $\mathcal{B}_+\oplus\mathcal{B}_-$ 对应正则表示，可以离开奇点形变，故有模。易错处：星号标记的是「在该一般点处支撑」，$\mathcal{B}^*$ 与 $\mathcal{B}$ 不是同一对象。]

若把 $\mathcal{B}$ 形变为支撑在邻近某条 Hitchin 纤维的光滑点 $r^*$ 上的点层，则 $\mathcal{A}$ 形变为支撑在对偶 Hitchin 纤维上的秩 $1$ $A$-膜，并带有一个由 $r^*$ 决定的平坦酉线丛。考察把 $\mathcal{B}$ 形变到同一奇异纤维的光滑点 $r^*$ 时会发生什么，也很有启发性。此时对偶的 $A$-膜是同一对偶（奇异）Hitchin 纤维上的平坦酉线丛。然而这个平坦线丛现在具有联系纤维的两个不可约分支的非平凡单值，因此它们不再能被「拆开」。换言之，奇异 Hitchin 纤维上一般的秩 $1$ $A$-膜实际上不可约，这与它对应于支撑在光滑点上的秩 $1$ $B$-膜（后者也不可约）这一事实相符。

\subsubsection{两个分支之间的关系}

作为超凯勒流形，两个分支 $\mathcal{M}_H(SO_3,0)$ 与 $\mathcal{M}_H(SO_3,\theta)$ 是不同的。但若仅把它们视为复结构 $I$ 下的复辛流形，则只要 $\sigma\ne0$，它们实际上同构。

可以通过在点 $p$ 处作一次 $\varphi$-不变的 Hecke 修改，在两个分支之间建立映射。这一概念在 §5.3.3 中有更充分的展开。简言之，把 $E$ 在 $p$ 附近分解为 $\mathcal{L}_1\oplus\mathcal{L}_2$，其中 $\mathcal{L}_1$ 与 $\mathcal{L}_2$ 是 $\varphi$-不变线丛，即 $\varphi:E\to E\otimes K$ 把 $\mathcal{L}_i$ 映入 $\mathcal{L}_i\otimes K$（$i=1,2$）。排列 $\mathcal{L}_i$ 的顺序，使 $\mathcal{L}_1$ 在 $p$ 处的纤维是 $\sigma$ 的属于特征值 $+\sigma_0$ 的特征子空间。（这一步需要 $\sigma\ne0$。）考虑这样的操作：在 $p$ 之外保持 $(E,\varphi)$ 不变，而在 $p$ 附近作用为 $\mathcal{L}_1\oplus\mathcal{L}_2\to\mathcal{L}_1(p)\oplus\mathcal{L}_2$。当视为 $SO_3$ Higgs 丛 $(W,\varphi)$（其中 $W=\mathrm{ad}(E)$）的变换时，这一操作建立了 $\mathcal{M}_H(SO_3)$ 两个分支之间的同构。该同构与 Hitchin 纤维化交换，因为 Hecke 修改不改变 $\varphi$ 的特征多项式。
% [存疑: p.28 此处写「$\sigma$ 的特征值 $+\sigma_0$ 的特征子空间」，按上下文 $\sigma$ 应指 $\varphi$ 在 $p$ 处的留数（其本征值为 $\pm\sigma_0$，见 §3.9.1 对留数特征空间的用法）；原文措辞如此，译文照录，不改符号。]

回到分歧 Higgs 丛的 Hitchin 方程，并按式 (3.4) 假定奇性，则 $\varphi$-不变的 Hecke 修改等价于把 $\alpha$ 平移一个格矢。（可以用在 $r=0$ 处不连续的规范变换补偿这种平移，该变换的效果是改变丛在该点处的自然扩张。见 [GW] 的 §2.1。）一般而言 [Nak]，只要 $\sigma$ 正则，把 $\mathcal{M}_H$ 视为
% ==== 原文 PDF 第 29 页（印刷页 29）====
复结构 $I$ 下的复辛流形时，它不依赖于 $\alpha$。$\mathcal{M}_H(SO_3)$ 两个分支的等价性是这一事实的特例。对非正则的 $\sigma$（对 $SO_3$ 而言即 $\sigma=0$），$\mathcal{M}_H(SO_3)$ 的两个分支双有理等价但不同构。

对非分歧的 $SO_3$ Higgs 丛，两个分支之间不存在这样的关系。

\subsection{\texorpdfstring{$O_2$}{O2}-丛}

下一个任务是更好地解释我们在 $SO_3$ 模空间中发现的奇点。

一般地，对一般的 $\sigma$，稳定分歧 Higgs 丛的模空间 $\mathcal{M}_H$ 的奇点完全来自自同构。\footnote{对非正则的 $\sigma$，$\mathcal{M}_H$ 还有 [GW] §3.6 所描述的局部奇点。} 若 Higgs 丛 $(E,\varphi)$ 具有非平凡有限自同构群 $\Gamma$，则应预期 $\mathcal{M}_H$ 中的对应点是奇点。若 $\Gamma$ 有限，则奇点是轨形奇点，局部形如 $\mathbb{C}^{2n}/\Gamma$（对某个 $n$ 以及 $\Gamma$ 在 $\mathbb{C}^{2n}$ 上的某个线性作用）。若 $\Gamma$ 的维数为正，则奇点通常（但并非总是）比轨形奇点更严重。
% [Tutor Note: 推导补足。$\Gamma$ 有限时的轨形奇点来自稳定化子群在法空间上的线性化作用：把 $\mathcal{M}_H$ 局部写成矩映射零集的辛商后，$\Gamma$ 的作用线性化即给出 $\mathbb{C}^{2n}/\Gamma$，此处 $2n=\dim_\mathbb{C}\mathcal{M}_H$。$A_1$ 奇点就是 $\mathbb{C}^2/\mathbb{Z}_2$（$\mathbb{Z}_2$ 以 $-1$ 作用），故一旦观测到 $A_1$ 奇点，就强制 $\Gamma$ 含 $\mathbb{Z}_2$；本题中 $O_2$ 的中心正是这个 $\mathbb{Z}_2$。]

在目前的情形，我们已经遇到一些 $A_1$ 轨形奇点，这强烈提示相应的 $SO_3$ Higgs 丛 $(W,\varphi)$ 的自同构群为 $\Gamma=\mathbb{Z}_2$。为描述具有这一自同构群的 Higgs 丛，我们使用由 Hitchin 方程给出的 Higgs 丛与局部系统之间的对应。该对应保持自同构群，因此自同构群为 $\mathbb{Z}_2$ 的 Higgs 丛对应于自同构群为 $\mathbb{Z}_2$ 的局部系统。

之所以能有自同构群为 $\mathbb{Z}_2$ 的 $SO_3$ 局部系统，是因为 $SO_3$ 含有子群 $O_2$，它由形如
\begin{equation}\tag{3.21}
\begin{pmatrix}
* & * & 0\\
* & * & 0\\
0 & 0 & \pm1
\end{pmatrix}
\end{equation}
的 $SO_3$ 元素组成，其中左上角方块是 $O_2$ 的元素，而右下角的符号选取得使行列式等于 $+1$。$O_2$ 在拓扑上有两个分支；与单位元连通的分支由右下角矩阵元为 $+1$ 的群元素组成，不连通分支则由该矩阵元为 $-1$ 的元素组成。$SO_3$ 中与 $O_2$ 交换的子群是 $\mathbb{Z}_2$，由
\begin{equation}\tag{3.22}
\begin{pmatrix}
-1 & 0 & 0\\
0 & -1 & 0\\
0 & 0 & 1
\end{pmatrix}
\end{equation}
生成。

因此，结构群约化到 $O_2$ 但在其他方面一般的 $SO_3$ 局部系统，其自同构群为 $\mathbb{Z}_2$。

下面描述亏格 $1$ 的 Riemann 面 $C$ 上的局部系统，先看非分歧情形。这样的局部系统在同构意义下由绕 $C$ 上两个一循环的单值元 $V_1$ 与 $V_2$ 决定。由于 $C$ 的基本群是
% ==== 原文 PDF 第 30 页（印刷页 30）====
交换的，它们满足
\begin{equation}\tag{3.23}
V_1V_2V_1^{-1}V_2^{-1}=1.
\end{equation}

现设只有一个分歧点 $p$，并指定绕 $p$ 的单值所在的共轭类 $M$。则 (3.23) 修改为
\begin{equation}\tag{3.24}
V_1V_2V_1^{-1}V_2^{-1}=M.
\end{equation}

在我们的情形，单值为 $M=\exp(-2\pi(\alpha-\mathrm{i}\gamma))$，如式 (3.5) 所示。
% [原文疑误: p.30 原文作 "as in in eqn. (3.5)"，重复了一个 "in"；译文按 "as in eqn. (3.5)" 处理。]
% [Tutor Note: 易错预警。(3.24) 中的 $M$ 只是绕分歧点单值的共轭类；一旦选定 $V_1,V_2$ 的具体表示，$M$ 由换位子 $V_1V_2V_1^{-1}V_2^{-1}$ 给出。若 $\mathcal{S}$ 平凡则 $V_1,V_2$ 都落在交换群 $SO_2$ 中，换位子必为 $1$，这就是「对 $M\ne1$ 不可能」的全部内容；反之 $\mathcal{S}$ 非平凡迫使某个 $V_i$ 落在 $O_2$ 的不连通分支，换位子才可能非平凡。]

若 $V_1$ 与 $V_2$ 取值于 $O_2$ 且满足 (3.24)，则 $M$ 必取值于 $O_2$ 的连通分支，并具有形式
\begin{equation}\tag{3.25}
M=
\begin{pmatrix}
* & * & 0\\
* & * & 0\\
0 & 0 & 1
\end{pmatrix}.
\end{equation}

对这样的 $M$，我们要找取值于 $O_2$ 且满足 (3.24) 的 $V_1$ 与 $V_2$。

结构群实际上约化到 $O_2$ 的 $SO_3$ 局部系统 $W$ 分解为 $U\oplus\mathcal{S}$，其中 $U$ 是秩二局部系统（结构群 $O_2\subset GL_2$），而 $\mathcal{S}\cong\det U$ 是 2 阶线丛。
% [原文疑误: p.30 原文作 "where $W$ is a rank two local system"，但 $W$ 是整个秩三的 $SO_3$ 局部系统，分裂出的秩二部分是 $U$；据上下文应为 "where $U$ is a rank two local system"，译文按 $U$ 处理。]
若 $\mathcal{S}$ 平凡，则 $V_1$ 与 $V_2$ 都取值于 $O_2$ 的连通分支。由于该连通分支是交换群 $SO_2$，此时对 $M\ne1$，(3.24) 不可能满足。

所以必须取 $\mathcal{S}$ 为非平凡的 2 阶线丛。这意味着 $V_1$ 或 $V_2$（或两者）取值于 $O_2$ 的不连通分支。我们可以选取和乐为 $V_1$、$V_2$ 的两个一循环，使 $V_1$ 取值于不连通分支而 $V_2$ 取值于连通分支。不连通分支的任一元素共轭于
\begin{equation}\tag{3.26}
V_1=
\begin{pmatrix}
1 & 0 & 0\\
0 & -1 & 0\\
0 & 0 & -1
\end{pmatrix}.
\end{equation}

取这样的 $V_1$ 时，有 $V_1V_2V_1^{-1}=V_2^{-1}$，于是 (3.24) 化为
\begin{equation}\tag{3.27}
V_2^2=M^{-1},
\end{equation}
其中 $M$ 与 $V_2$ 取值于 $SO_2$。

对每个 $M$，最后一个方程恰有两个解，它们在 $SO_2$ 中相差 $V_2\to-V_2$。因此，对每个非平凡的 2 阶线丛 $\mathcal{S}$，我们可以构造恰两个 $SO_3$ 局部系统，各自具有自同构群 $R\cong\mathbb{Z}_2$。

然而这两个局部系统的 $w_2$ 值不同。因此，其中之一对应于 $\mathcal{M}_H(C;SO_3,0)$ 相应纤维上的一个 $\mathbb{Z}_2$ 轨形奇点，另一个则对应于 $\mathcal{M}_H(C;SO_3,\theta)$ 相应纤维上的此类奇点。

为说明为何 $V$ 的两种选择对应拓扑上不等价的丛，考虑特殊情形 $M=1$。此时 $V_2$ 的两种可能为
\begin{equation}\tag{3.28}
V_2^{(1)}=1,\qquad
V_2^{(2)}=
\begin{pmatrix}
-1 & 0 & 0\\
0 & -1 & 0\\
0 & 0 & 1
\end{pmatrix}.
\end{equation}

% ==== 原文 PDF 第 31 页（印刷页 31）====

两种情形下 $V_1$ 与 $V_2$ 都是对角的，故 $W$ 分裂为三个秩 $1$ 局部系统 $\mathcal{S}_1$、$\mathcal{S}_2$、$\mathcal{S}_3$ 的直和。于是有
\begin{equation}\tag{3.29}
w_2(W)=\sum_{1\le i<j\le3}w_1(\mathcal{S}_i)\cup w_1(\mathcal{S}_j).
\end{equation}
% [Tutor Note: 推导补足。(3.29) 是正交丛分裂为线丛时的 Stiefel--Whitney 公式：$w(W)=\prod_i(1+w_1(\mathcal{S}_i))$，取 2 次项即得 $\sum_{i<j}w_1(\mathcal{S}_i)\cup w_1(\mathcal{S}_j)$。易错：$\cup$ 是 $\mathbb{Z}_2$ 系数上同调的杯积，故 $H^1$ 中元素满足 $x^2=0$，这正是 $V_2=1$ 时 $w_2(W)=w_1(\mathcal{S}_2)^2=0$ 的原因；只有 $x,y$ 构成 $H^1(C,\mathbb{Z}_2)$ 的基时，$\theta=x\cup y$ 才非零。]

若 $V_2=1$，则 $\mathcal{S}_1$ 平凡，且 $\mathcal{S}_2$ 与 $\mathcal{S}_3$ 同构。于是 $w_2(W)=w_1(\mathcal{S}_2)^2$，而对任意 $x\in H^1(C,\mathbb{Z}_2)$ 这为零，因为 $x^2=0$。

若 $V_2=V_2^{(2)}$，则 $\mathcal{S}_1$、$\mathcal{S}_2$、$\mathcal{S}_3$ 是 $C$ 上三个互不相同的非平凡 2 阶线丛。令 $x=w_1(\mathcal{S}_1)$、$y=w_1(\mathcal{S}_2)$，则由 $\mathcal{S}_3\cong\mathcal{S}_1\otimes\mathcal{S}_2$ 得 $w_1(\mathcal{S}_3)=x+y$。$x$ 与 $y$ 构成 $H^1(C,\mathbb{Z}_2)\cong\mathbb{Z}_2\oplus\mathbb{Z}_2$ 的一组基，而 $\theta=x\cup y$ 生成 $H^2(C,\mathbb{Z}_2)\cong\mathbb{Z}_2$。代入 (3.29) 得 $w_2(W)=x\cup y$，特别地 $w_2(W)\ne0$。

于是，对任意非平凡 2 阶线丛 $\mathcal{S}$ 的选择以及任意 $w_2(W)$ 的选择，我们都得到一个自同构群为 $\mathbb{Z}_2$ 的 $SO_3$ 局部系统。这解释了我们在 §3.6 对 $w_2=0$ 找到的奇点，以及我们将对 $w_2\ne0$ 找到的奇点。

\par\medskip
\noindent\textbf{注 3.2.}\ \textit{在 $w_2\ne0$ 的例子中，我们可以不把 $W=\mathcal{S}_1\oplus\mathcal{S}_2\oplus\mathcal{S}_3$ 看作 $SO_3$ 局部系统，而把它理解为其结构群为 $SO_3$ 的紧实形式的平坦丛，或等价地理解为稳定的全纯 $SO_3$ 丛。如此，$W$ 是亏格 $1$ 曲线上唯一的拓扑非平凡平坦丛或稳定全纯 $SO_3$ 丛。特别地，它没有形变；例如 $H^1(C,W)=\oplus_{i=1}^3H^1(C,\mathcal{S}_i)=0$，因为对亏格 $1$ 的 $C$，非平凡线丛的上同调为零。}
\par\medskip

\par\medskip
\noindent\textbf{注 3.3.}\ \textit{我们在描述 $\mathcal{M}_H(C;SO_3)$ 两个分支的区分时过于随意了。在非分歧情形，$w_2(W)$ 是定义良好的拓扑不变量。但在分歧情形，$w_2(W)$ 的定义依赖于 $W$ 在分歧点 $p$ 上方扩张（至少作为拓扑丛的扩张）的选取。模空间 $\mathcal{M}_H(C;SO_3)$ 有两个分支，每个分支都具有我们所描述的奇点，但断言其中一个对应 $w_2(W)=0$、另一个对应 $w_2(W)\ne0$ 则过于简化。事实上，两个分支在 $M$ 于群 $SO_2$ 中的单值下互换，而 $SO_2$ 的基本群为 $\mathbb{Z}$。$M$ 的单值可通过改变参数 $\alpha$ 实现，因此等价于 §3.6.2 所描述的操作。}
\par\medskip

\subsubsection{与几何内窥的关系}

我们所发现的，再次与 §2.2 的预期完全一致。结构群约化到 $O_2$ 的 $SO_3$ 局部系统具有自同构群 $\mathbb{Z}_2$。它对应于 $\mathcal{M}_H(C;SO_3)$ 的一个 $A_1$ 奇点。因此与该局部系统相联系的 $B$-膜范畴由两个对象 $\mathcal{B}_+$ 与 $\mathcal{B}_-$ 生成。相比之下，$\mathcal{M}_H(C;SL_2)$ 光滑；Hitchin 纤维化的某些例外纤维奇异，但纤维的奇点是全空间的光滑点。

\subsection{第二分支}

在 §3.7 中，我们构造了 $w_2\ne0$、结构群约化到 $O_2$ 的 $SO_3$ 局部系统。它们应当对应于 $\mathcal{M}_H(C;SO_3,\theta)$ 的 $A_1$ 奇点。

% ==== 原文 PDF 第 32 页（印刷页 32）====

给出这个空间的直接描述将是有用的。$\mathcal{M}_H(C;SO_3,\theta)$ 参数化分歧 Higgs 丛 $(W,\varphi)$，其中 $W$ 是 $w_2\ne0$ 的 $SO_3$-丛。对我们的目的而言，只需在该模空间中构造一个 Zariski 开集，它对应于 $W$ 稳定的那些 Higgs 丛。据注 3.2，这意味着 $W=\mathcal{S}_1\oplus\mathcal{S}_2\oplus\mathcal{S}_3$ 是三个非平凡 2 阶线丛的直和。可以取 $\mathcal{S}_i=\mathcal{O}(p)\otimes\mathcal{O}(q_i)^{-1}$，其中 $q_i$ 是由 $(x,y)=(e_i,0)$ 给出的 2 阶点。作为 $SO_3$ 丛，$W$ 应具有非退化二次型，我们把它记作迹（即把 $W$ 视为伴随丛）。其定义如下：若 $s=s_1\oplus s_2\oplus s_3$ 是 $W=\mathcal{S}_1\oplus\mathcal{S}_2\oplus\mathcal{S}_3$ 的截面，则二次型为 $\Tr s^2=\sum_i s_i^2/(x-e_i)$。其想法是：$s_i^2$ 是 $\mathcal{O}(p)^2\otimes\mathcal{O}(q_i)^{-2}$ 的截面，而除以 $x-e_i$ 即可将其平凡化。
% [Tutor Note: 推导补足。$\Tr s^2=\sum_i s_i^2/(x-e_i)$ 不是朴素平方和：$\mathcal{S}_i=\mathcal{O}(p)\otimes\mathcal{O}(q_i)^{-1}$，故 $s_i^2$ 是 $\mathcal{O}(p)^2\otimes\mathcal{O}(q_i)^{-2}$ 的截面，在 $q_i$ 处有零点；除以在 $q_i$ 处单零的 $x-e_i$ 后才成为 $\mathcal{O}(p)^2\cong K^2$ 的截面，这正与 $\varphi\in H^0(C,K\otimes\mathcal{O}(p)\otimes W)$ 所配的二次型一致。易错：$q_i$ 是 $(x,y)=(e_i,0)$ 给出的 2 阶点，不要与 $e_i$ 本身混同。]

丛 $W$ 具有自同构群 $Q=\mathbb{Z}_2\times\mathbb{Z}_2$。事实上，作为平坦丛，$W$ 具有式 (3.26) 与 (3.28) 给出的单值 $V_1$ 与 $V_2^{(2)}$，而 $Q$ 是由对角元为 $\pm1$ 的对角矩阵组成的 $SO_3$ 子群。

为构造 $\mathcal{M}_H(C;SO_3,\theta)$（更确切地说，对应于稳定丛 $E$ 的 Zariski 开集），我们必须考虑所有 Higgs 场 $\varphi\in H^0(C,K\otimes\mathcal{O}(p)\otimes W)$，对 $\Tr\varphi^2$ 的极点部分施加适当条件，再除以 $Q$。在着手这样做之前，先讨论不除以 $Q$ 时会发生什么。

从复代数几何的观点看，与其考虑 $E$ 为行列式平凡的秩二全纯丛的 $SL_2$ Higgs 丛 $(E,\varphi)$，不如取定线丛 $\mathcal{L}$ 并考虑 $\det E=\mathcal{L}$ 的 Higgs 丛 $(E,\varphi)$。若把 $\mathcal{L}$ 换成 $\mathcal{L}\otimes\mathcal{N}^2$，可通过 $E\to E\otimes\mathcal{N}$ 补偿。模去这一操作，$\mathcal{L}$ 真正重要的只是其次数模 $2$。除熟悉的 $\mathcal{L}=\mathcal{O}$ 情形外，我们还可以考虑第二个分支，比如取 $\det E=\mathcal{O}(p)$。我们把这类稳定 Higgs 丛的模空间记作 $\mathcal{M}_H(C;SL_2^*)$。有时称之为 $SL_2$ 模空间的 improper 分支（improper component）（而把 $\det E$ 平凡的通常 $SL_2$ Higgs 丛 $(E,\varphi)$ 的模空间称为 proper 分支）。若 $(E,\varphi)$ 是 $\det E=\mathcal{O}(p)$ 的 Higgs 丛，则令 $W=\mathrm{ad}(E)$ 便得到 $w_2(W)\ne0$ 的 $SO_3$ Higgs 丛 $(W,\varphi)$。所有 $(W,\varphi)$ 都可由此得到，而 $(E,\varphi)$ 在相差把 $E$ 扭转一个 2 阶线丛的意义下由 $(W,\varphi)$ 唯一确定。因此，通过与 (3.20) 类比，可把 $\mathcal{M}_H(C;SL_2^*)$ 与 $SO_3$ Higgs 丛联系起来：
\begin{equation}\tag{3.30}
\mathcal{M}_H(C;SO_3,\theta)=\mathcal{M}_H(C;SL_2^*)/Q.
\end{equation}

从微分几何的观点看，$\mathcal{M}_H(C;SL_2^*)$ 可通过在 $w_2\ne0$ 的丛上以紧群 $SO_3$ 为规范群求解 Hitchin 方程、并只对可提升到 $SU(2)$ 的规范变换取商来构造。

\par\medskip
\noindent\textbf{注 3.4.}\ \textit{注 3.3 的论证也可施于群 $SL_2$，以说明 $\mathcal{M}_H(C;SL_2)$ 与 $\mathcal{M}_H(C;SL_2^*)$ 作为复结构 $I$ 下的复辛流形同构。}
\par\medskip

现在给出显式描述。考虑作为 $K\otimes\mathcal{O}(p)\otimes W=\oplus_{i=1}^3K\otimes\mathcal{O}(p)^2\otimes\mathcal{O}(q_i)^{-1}$ 的截面的 Higgs 场 $\varphi$。对每个 $i$，可取 $K\otimes\mathcal{O}(p)^2\otimes\mathcal{O}(q_i)^{-1}$ 的一个截面 $u_i$，即 $u_i=(dx/y)(x-e_i)$，并有 $\Tr u_i^2=(dx/y)^2(x-e_i)$。Higgs 场的一般形式为
% ==== 原文 PDF 第 33 页（印刷页 33）====
\begin{equation}\tag{3.31}
\varphi=\sum_{i=1}^{3}a_iu_i,
\end{equation}
其中 $a_i$ 为复常数。这给出
\begin{equation}\tag{3.32}
\Tr\varphi^2=\left(\frac{dx}{y}\right)^2\sum_{i=1}^{3}a_i^2(x-e_i).
\end{equation}

如式 (3.11) 那样，乘以 $(dx/y)^2$ 的是 $x$ 的一次多项式。令 $z\sim x^{-1/2}$ 为无穷远处的局部参数，则 $\Tr\varphi^2$ 的极点部分为 $4(dz/z)^2\sum_i a_i^2$。令其等于 $2\sigma_0^2(dz/z)^2$，便要求
\begin{equation}\tag{3.33}
\sum_i a_i^2=\frac{\sigma_0^2}{2}.
\end{equation}

这个仿射二次曲面描述了 $\mathcal{M}_H(C;SL_2^*)$ 中的一个 Zariski 开集（在 §3.8.1 中，我们将除以 $Q=\mathbb{Z}_2\times\mathbb{Z}_2$，并给出 $\mathcal{M}_H(C;SO_3,\theta)$ 的类似描述）。

(3.32) 右端乘以 $(dx/y)^2$ 的常数项是 $-\sum_ie_ia_i^2$。这使我们能描述 Hitchin 纤维化；它是由 $(a_1,a_2,a_3)$ 到
\begin{equation}\tag{3.34}
w_0=-\sum_i e_ia_i^2
\end{equation}
的映射。

Hitchin 纤维化的纤维由两个二次曲面 (3.33) 与 (3.34) 的交给出。出于讨论 $SL_2$ 时的同样理由，假定 $\sigma_0$ 非零，则可把这些方程中的参数 $\sigma_0$ 消去。令 $w_0=-\sigma_0^2w/2$，$b_i=a_i(\sqrt2/\sigma_0)$，即把两个二次曲面写成
\begin{equation}\tag{3.35}
\begin{aligned}
b_1^2+b_2^2+b_3^2&=1\\
e_1b_1^2+e_2b_2^2+e_3b_3^2&=w.
\end{aligned}
\end{equation}
% [Tutor Note: 易错预警。由 (3.33) 的 $\sum a_i^2=\sigma_0^2/2$ 得 $\sum b_i^2=\sum a_i^2(2/\sigma_0^2)=1$；由 (3.34) 的 $w_0=-\sum e_ia_i^2=-(\sigma_0^2/2)\sum e_ib_i^2$，故 $w_0=-\sigma_0^2w/2$ 与 $\sum e_ib_i^2=w$ 一致。符号、因子 $2$ 与 $\sqrt2$ 必须联立核对；另外 $z\sim x^{-1/2}$ 来自 $K\otimes\mathcal{O}(p)$ 的谱覆盖上无穷远处的局部参数满足 $z^2\sim x^{-1}$。]

对一般的 $w$，该交是一条亏格 $1$ 的光滑曲线，其中有些点被略去，因为我们假定了 $W$ 稳定。$w$ 取什么值时纤维奇异？若 $f=\sum_ib_i^2-1$、$g=\sum_ie_ib_i^2-w$，则纤维的奇点是满足 $f=g=df\wedge dg=0$ 的点。简短计算表明，$df\wedge dg=0$ 当且仅当 $b_1$、$b_2$、$b_3$ 中有两个为零。若对某个 $i$，$b_i$ 非零而对 $j\ne i$ 有 $b_j=0$，则必须有
\begin{equation}\tag{3.36}
w=e_i
\end{equation}
以及
\begin{equation}\tag{3.37}
b_i=\pm1.
\end{equation}

于是恰有三条奇异纤维，每个 $i$ 对应一条，与 §3.4 相同。奇异纤维位于 $w=e_i$（$i=1,2,3$），与我们为模空间的另一分支所找到的一致。这些事实当然与 §3.6.2 所断言的两分支之间的关系相符。

此外，每条奇异纤维 $\boldsymbol{F}_w$ 含有两个奇点，由式 (3.37) 给出。奇异纤维由两个亏格 $0$ 分支组成，它们通过两个二重点相接。

% ==== 原文 PDF 第 34 页（印刷页 34）====

为看出这一点，取 $i=1$。若 $w=e_1$，则方程 $f=0$ 与 $g=0$ 的某个线性组合给出 $(e_2-e_1)b_2^2+(e_3-e_1)b_3^2=0$，即
\begin{equation}\tag{3.38}
b_2=\pm b_3\sqrt{-(e_3-e_1)/(e_2-e_1)}.
\end{equation}
这描述了一条曲线 $\boldsymbol{F}_{w,0}$，它是两个亏格零分支在一个点 $b_2=b_3=0$ 处相交的并。再由 $b_1^2=1-b_3^2(e_2-e_3)/(e_2-e_1)$ 解出 $b_1$，便得到 $\boldsymbol{F}_{w,0}$ 的一个双重覆盖。该双重覆盖（即 Hitchin 纤维化的纤维 $\boldsymbol{F}_w$）在 $\boldsymbol{F}_{w,0}$ 每个分支上的两个点处分歧。亏格零曲线在两点处分歧的双重覆盖仍为亏格零。因此 $\boldsymbol{F}_w$ 由两个亏格零分支组成，它们在两个点 $b_2=b_3=0$、$b_1=\pm1$ 处相接。

\subsubsection{\texorpdfstring{$SO_3$}{SO3} 模空间}

迄今我们构造了 $\mathcal{M}_H(C;SL_2^*)$ 中的一个 Zariski 开集及其 Hitchin 纤维化。我们还可以除以 $Q$——它通过成对改变 $b_i$ 的符号作用——从而得到 $\mathcal{M}_H(C;SO_3,\theta)$ 中的一个 Zariski 开集。

在成对改变 $b_i$ 符号下的不变量是 $b_1^2$、$b_2^2$、$b_3^2$ 以及 $z=b_1b_2b_3$。它们满足 $z^2=b_1^2b_2^2b_3^2$。利用两个方程 (3.35)，可把 $b_i^2$（$i=1,2,3$）解为 $t=\sum_{i=1}^3e_i^2b_i^2$ 的线性函数。方程 $z^2=b_1^2b_2^2b_3^2$ 化为
\begin{equation}\tag{3.39}
z^2=-\prod_{1\le i<j\le3}\frac{(t-w(e_i+e_j)+e_ie_j)}{(e_i-e_j)^2}.
\end{equation}

这描述了 $SO_3$ 模空间（只是还须计入固定 $w$ 时 $t,z=\infty$ 的点）。

用这一描述来找出 $\mathcal{M}_H(C;SO_3,\theta)$ 的 Hitchin 纤维化的奇异纤维。(3.39) 右端的三次多项式除非对某个 $i$ 有 $w=e_i$（此时恰有两个根重合），否则根互不相同。因此唯一的奇异纤维出现在 $w=e_i$ 处，它们是节点三次曲线，即把 $\mathbb{C}\mathbb{P}^1$ 的两个点等同所得的曲线。在二重点附近展开三次方程 (3.39)，不难验证纤维的二重点实际上是 $\mathcal{M}_H(C;SO_3,\theta)$ 的 $A_1$ 奇点。这与我们在 §3.6 末为 $\mathcal{M}_H(C;SO_3,0)$ 的奇异纤维所得的结果相同。这当然与 §3.6.2 关于模空间两分支在复结构 $I$ 下等价的断言完全相符。

\subsection{与余切丛的关系}

稳定 Higgs 丛（无分歧）的模空间 $\mathcal{M}_H$ 的最重要性质之一 [Hi1] 是它可以近似为 $T^*\mathcal{M}$，其中 $\mathcal{M}$ 是稳定丛的模空间。（$\mathcal{M}$ 是参数化稳定与半稳定丛的模空间中的一个 Zariski 开集。）
% [原文疑误: p.34 原文作 "a Zariski open set the moduli space that parametrizes ..."，脱漏介词；译文按 "a Zariski open set in the moduli space" 理解处理。]
理由是：与稳定丛 $E$ 对应的点处 $\mathcal{M}$ 的余切空间为 $H^0(C,K\otimes\mathrm{ad}(E))$。于是 $T^*\mathcal{M}$ 中的点是一对 $(E,\varphi)$，换言之即一个 Higgs 丛。这给出 $T^*\mathcal{M}$ 作为 $\mathcal{M}_H$ 中 Zariski 开集的嵌入。该映射不满，因为即使底丛 $E$ 不稳定，Higgs 丛 $(E,\varphi)$ 也可能稳定或半稳定，从而代表 $\mathcal{M}_H$ 中的点。

分歧 Higgs 丛有类似的对应。此时取 $\mathcal{M}$ 为在一点 $p$ 处带抛物结构的稳定丛的模空间，$\mathcal{M}_H$ 为稳定分歧 Higgs 丛的模空间。（对一般的抛物权选取——即式 (3.4) 中称为 $\alpha$ 的参数——每个半稳定抛物丛都是稳定的。）
% ==== 原文 PDF 第 35 页（印刷页 35）====
于是 $\mathcal{M}_H$ 有一个 Zariski 开集，它不是 $T^*\mathcal{M}$，而是 $T^*\mathcal{M}$ 的仿射辛形变，如 [GW] 的 §3.6 所述。我们把这样的仿射辛形变记作 $\tilde{T}^*\mathcal{M}$。这里 $\tilde{T}^*\mathcal{M}$ 表示一个带映射到 $\mathcal{M}$ 的复辛簇，使得在 $\mathcal{M}$ 的局部它辛同构于 $T^*\mathcal{M}$。关于一个例子的详细描述，见 §3.9.1。
% [背景补充: 「仿射辛形变」指这样的复辛簇：它到 $\mathcal{M}$ 有映射，且在 $\mathcal{M}$ 的每个 Zariski 开集上辛同构于余切丛 $T^*\mathcal{M}$，但整体上只相差一个仿射线性变换（本例中由 $\mathbb{C}$ 的平移实现）。它与通常的形变不同：形变改变辛结构，仿射形变只改变余切丛的零截面。]

对几何 Langlands 的应用而言，重要的是把 Hitchin 纤维化的纤维从 $\mathcal{M}_H$ 限制到 $T^*\mathcal{M}$ 或 $\tilde{T}^*\mathcal{M}$，因为这是用 D-模解释 A-膜的关键一步，如 §2.3 所述。下面就对我们这个例子执行这一步：亏格 $1$ 曲线 $C$ 上带一个分歧点 $p$ 的 $SL_2$ Higgs 丛。当然，也可以除以 $Q$ 并对 $SO_3$ 作类似讨论。
% [原文疑误: p.35 原文作 "as discussed in Section 2.3)."，多出一个右括号；译文按 "as discussed in Section 2.3" 处理。]

\subsubsection{计算}

这些问题只依赖于 $\mathcal{M}_H$ 作为复结构 $I$ 下的复辛流形。按 §3.6.2，两个分支在此意义下等价，因此从现在起我们考虑 improper 分支。

我们已显式描述了 $\det E=\mathcal{O}(p)$ 且 $E$ 稳定的稳定 Higgs 丛 $(E,\varphi)$。$E$ 在同构意义下唯一确定，而我们已经发现这类 Higgs 丛构成的空间恰是仿射二次曲面
\begin{equation}\tag{3.40}
b_1^2+b_2^2+b_3^2=1,
\end{equation}
暂时把它记作 $\mathcal{M}_H^0$。它与 $\mathcal{M}_H$ 略有不同：后者还参数化底丛 $E$ 不稳定的稳定对 $(E,\varphi)$。

在此情形，稳定抛物丛的模空间 $\mathcal{M}$ 是什么？丛 $E$ 没有形变，但其在点 $p$ 处的抛物结构可以变动。抛物结构的选取即是在二维向量空间 $E_p$（$E$ 在 $p$ 处的纤维）中选取一条复直线 $\Psi$。所以 $\mathcal{M}$ 是一份 $\mathbb{C}\mathbb{P}^1$，参数化 $\Psi$ 的选取。

取如上之 $E$，设给定分歧 Higgs 丛 $(E,\varphi)$，其中 $\varphi$ 在 $p$ 处有极点，其留数的特征值为 $\pm\sigma_0$。这就在 $E$ 上确定一个抛物结构：取 $\Psi$ 为留数属于特征值 $+\sigma_0$ 的特征子空间。于是有自然映射 $\mathcal{M}_H^0\to\mathcal{M}$，把 $(E,\varphi)$ 映为以此 $\Psi$ 为抛物结构的抛物丛。该映射没有全纯截面，因为给定丛 $E$ 与 $\Psi$ 的选取，并没有自然的方式构造出具有正确极点的 $\varphi$。正是这一点导致了仿射形变。

或者，用属于特征值 $-\sigma_0$ 的特征子空间，可以得到第二个同样自然的抛物结构与映射 $\mathcal{M}_H^0\to\mathcal{M}$。这最终会导向从 A-膜到扭 D-模的不同函子。一般地，在几何 Langlands 的分歧情形，B-模型与 A-模型之间的镜像对称是唯一确定的，但从 A-膜到扭 D-模的映射依赖于包含局部单值的 Borel 子群的选取。不同的选取导致具有不同扭转的 D-模，它们之间由 Weyl 变换相联系。

在目前语境下，Borel 子群的选取归结为特征值 $+\sigma_0$ 或 $-\sigma_0$ 的选取。（$\mathcal{M}_H$ 两分支之间的同构也依赖于这一符号选择。）对 $SL_n$，类似的做法是选取 $\varphi$ 极点部分各特征值的一个排序。

% ==== 原文 PDF 第 36 页（印刷页 36）====

从仿射二次曲面 $b_1^2+b_2^2+b_3^2=1$ 到 $\mathbb{C}\mathbb{P}^1$ 的映射当然可以不用 Higgs 丛来描述。方程 $b_1^2+b_2^2+b_3^2=1$ 等价于说矩阵
\begin{equation}\tag{3.41}
M=
\begin{pmatrix}
b_1 & b_2-\mathrm{i}b_3\\
b_2+\mathrm{i}b_3 & -b_1
\end{pmatrix}
\end{equation}
的特征值为 $\pm1$。取 $\Psi$ 为方程 $M\psi=\psi$ 的解空间，便得到从仿射二次曲面到 $\mathbb{C}\mathbb{P}^1$ 的映射。若令 $\psi=\begin{pmatrix}1\\ z\end{pmatrix}$，则发现在仿射坐标下，到 $\mathbb{C}\mathbb{P}^1$ 的映射可由 $z=(1-b_1)/(b_2-\mathrm{i}b_3)=(b_2+\mathrm{i}b_3)/(1+b_1)$ 描述。
% [Tutor Note: 易错预警。(3.41) 中的 $M$ 是 $2\times2$ 矩阵，与花体的模空间 $\mathcal{M}$（及 $\mathcal{M}_H$）是不同记号，原文在此重用 $M$ 容易误读。取 $\psi=(1,z)^{\mathsf T}$ 把 $M\psi=\psi$ 的解空间实现为 $\mathbb{C}\mathbb{P}^1$ 上的点，$z$ 即仿射坐标；在 $b_1=-1$ 处分母为零，该坐标失效，故只能在 Zariski 开集上描述。]

从 $\mathcal{M}_H^0$ 到 $\mathbb{C}\mathbb{P}^1$ 的映射的纤维是一份份 $\mathbb{C}$。例如 $z=0$ 的纤维由 $b_1=1$、$b_2+\mathrm{i}b_3=0$ 给出，并由 $b_2$ 参数化。然而 $\mathcal{M}_H^0$ 不是 $\mathbb{C}\mathbb{P}^1$ 上的全纯线丛，因为映射 $\mathcal{M}_H^0\to\mathbb{C}\mathbb{P}^1$ 没有全纯截面。相反，丛 $\mathcal{M}_H^0\to\mathbb{C}\mathbb{P}^1$ 的结构群是 $\mathbb{C}$ 的仿射线性运动群，即形如 $x\to ax+b$（$a\in\mathbb{C}^\times$，$b\in\mathbb{C}$）的变换群。我们可以把 $\mathcal{M}_H^0$ 描述为秩 $1$ 的仿射丛。

仿射线性运动群有到 $\mathbb{C}^\times$ 的同态（忘掉 $b$）。因此，对秩 $1$ 仿射丛，总有一个典范相伴的、结构群为 $\mathbb{C}^\times$ 的纤维丛。在二次曲面的情形，相伴的簇是 $\mathbb{C}\mathbb{P}^1$ 的余切丛。二次曲面本身带有全纯辛形式：\footnote{乘性常数 $\sigma_0$ 是通过在由 $b_1,b_2,b_3$ 为实数所刻画的二循环上积分 $\Omega$，并与 [GW] 式 (3.77) 中 $\Omega$ 的上同调类的结果比较而确定的。}
\begin{equation}\tag{3.42}
\Omega=\mathrm{i}\sigma_0\frac{db_1\wedge db_2}{b_3}
\end{equation}
它是模空间上超凯勒结构的一部分。在 $\mathbb{C}\mathbb{P}^1$ 局部可以选取纤维化 $\mathcal{M}_H^0\to\mathbb{C}\mathbb{P}^1$ 的一个截面，并把 $\mathcal{M}_H^0$ 与 $\mathbb{C}\mathbb{P}^1$ 的余切丛等同。

为在 $b_1\ne-1$ 的 Zariski 开集（在此 $z=(1-b_1)/(b_2-\mathrm{i}b_3)=(b_2+\mathrm{i}b_3)/(1+b_1)$ 有定义）上显式地做到这一点，令 $v=-\sigma_0(b_2-\mathrm{i}b_3)$。小计算表明 $\Omega=dz\wedge dv$，于是在 $b_1=-1$ 之外，可把 $\mathbb{C}\mathbb{P}^1$ 作为轨迹 $v=0$ 嵌入 $\mathcal{M}_H^0$，并把 $\mathcal{M}_H^0$ 与 $z$-平面的余切丛等同。这一等同不能延拓\footnote{如上一条脚注所指出的，$\Omega$ 具有非零周期，这意味着它不能整体地成为某个余切丛的辛形式。}到 $z=\infty$ 之上；整体地，$\mathcal{M}_H^0$ 就是我们所称的余切丛的仿射形变。方便起见，再引入 $\tilde v=v/\sigma_0=-(b_2-\mathrm{i}b_3)$。

由于 $\mathcal{M}_H^0$ 不是 $\mathcal{M}$ 的余切丛，而是其仿射形变，此后我们把它记作 $\tilde{T}^*\mathcal{M}$。

现在考虑 Hitchin 纤维化的纤维。它们与二次曲面 $\tilde{T}^*\mathcal{M}$ 的交由在二次曲面的定义方程之外再补上方程
\begin{equation}\tag{3.43}
e_1b_1^2+e_2b_2^2+e_3b_3^2=w,
\end{equation}
% ==== 原文 PDF 第 37 页（印刷页 37）====
得到，从而给出代数曲线 $\boldsymbol{F}_w$。这条曲线当然可以投影到 $\mathcal{M}$，给出 $\mathcal{M}$ 的双重覆盖。然而一般而言，Hitchin 纤维化的纤维是完备（且一般光滑）的代数簇。这是 [Hi1] 理论的重要部分。我们将在目前语境下于 §5.3 用谱覆盖来解释它。

因此对一般的 $w$，$\mathcal{M}_H$ 的 Hitchin 纤维化的纤维 $\boldsymbol{F}_w$ 实际上是刚才描述的仿射曲线所对应的光滑射影曲线。用仿射二次曲面描述时会缺掉几个点。

可以通过再添加一个变量 $b_4$ 把 $\boldsymbol{F}_w$ 描述为射影曲线，其中 $b_1,\dots,b_4$ 理解为 $\mathbb{C}\mathbb{P}^3$ 上的齐次坐标并满足
\begin{equation}\tag{3.44}
\begin{aligned}
\sum_{i=1}^3b_i^2&=b_4^2\\
\sum_{i=1}^3e_ib_i^2&=wb_4^2.
\end{aligned}
\end{equation}

当用余切丛的仿射形变近似 $\mathcal{M}_H$ 时，缺失的是 $b_4=0$ 的四个点。这些点对应于抛物丛 $E$ 不稳定的稳定 Higgs 丛 $(E,\varphi)$。\footnote{$E$ 的行列式必须为 $\mathcal{O}(p)$，且必须形如 $\mathcal{L}\otimes(\mathcal{O}(p)\oplus\mathcal{O})$，其中 $\mathcal{L}$ 为 2 阶，才能存在稳定分歧 Higgs 丛 $(E,\varphi)$。$\mathcal{L}$ 的四种选择导致 Hitchin 纤维化每条纤维上的四个缺失点。} 它们构成 $b_1,b_2,b_3$ 成对变号群 $Q=\mathbb{Z}_2\times\mathbb{Z}_2$ 的单一轨道。具体地，这四个点对应的 $z$ 值为
\begin{equation}\tag{3.45}
z=\pm\sqrt{\frac{e_2-e_1}{e_2-e_3}}\pm\sqrt{\frac{e_3-e_1}{e_2-e_3}}.
\end{equation}

为了对几何 Langlands 纲领中出现的 D-模有所洞察，我们必须把 $\boldsymbol{F}_w$ 描述为 $\tilde{T}^*\mathcal{M}$ 中的曲线。为此使用坐标 $z$、$\tilde v$；经过一些代数运算发现，可以用关于 $\tilde v$ 的一个显式二次方程来刻画纤维 $\boldsymbol{F}_w$，其形式为
\begin{equation}\tag{3.46}
A(z)\tilde v^2+B(z)\tilde v+C(z)=0,
\end{equation}
其中
\begin{equation}\tag{3.47}
\begin{aligned}
A(z)&=(e_2-e_3)z^4+(4e_1-2e_2-2e_3)z^2+(e_2-e_3)\\
B(z)&=-4z((e_2-e_3)z^2+2e_1-e_2-e_3)\\
C(z)&=4((e_2-e_3)z^2+e_1-w).
\end{aligned}
\end{equation}
注意
\begin{equation}\tag{3.48}
B=-\frac{dA}{dz}.
\end{equation}
% [Tutor Note: 易错预警。(3.47) 中的 $B(z)=-4z((e_2-e_3)z^2+2e_1-e_2-e_3)$ 与 (3.48) $B=-dA/dz$ 是自洽性检验：对 $A(z)=(e_2-e_3)z^4+(4e_1-2e_2-2e_3)z^2+(e_2-e_3)$ 求导即得。四个缺失点恰为 $A$ 的零点，在那些点处 $B\ne0$，故二次方程一根有限、一根有极点，这正是 (3.49) 的来源。]

若令 $z$ 趋近式 (3.45) 中四个值之一，则关于 $\tilde v$ 的二次方程有一个根趋于无穷。因此在那些 $z$ 值处，若限制到 $\tilde{T}^*\mathcal{M}$，Hitchin 纤维中有一个点「缺失」。事实上，$z$ 的这四个临界值恰是多项式 $A(z)$ 的零点。在这些 $z$ 值处 $B(z)$ 非零，
% ==== 原文 PDF 第 38 页（印刷页 38）====
这意味着关于 $\tilde v$ 的二次方程 (3.46) 的两个根中一个有一阶极点，另一个没有。方便的是用变量 $v=\sigma_0\tilde v$ 来表达结果，它把辛形式 $\Omega$ 化为标准形式。设 $z^*$ 为 $A(z)$ 的任一零点。$v$ 的极点分支在 $z^*$ 附近的行为为 $v\sim-\sigma_0B(z)/A(z)$，利用式 (3.48) 可化为
\begin{equation}\tag{3.49}
v\sim\sigma_0\frac{1}{z-z^*}.
\end{equation}
% [Tutor Note: 物理图像。(3.49) 说明 Hitchin 纤维限制到 $\tilde{T}^*\mathcal{M}$ 时，在缺失除子（本例中为点 $z=z^*$）附近以 $v\sim\sigma_0/(z-z^*)$ 发散；这个一阶极点的留数 $\sigma_0$ 就是 §4 中扭 D-模的扭转参数。A-膜与 D-模的对应正是把 Hitchin 纤维的「无穷远行为」翻译成 D-模的正则奇点，这是 §4 全篇的技术起点。]

\par\medskip
\noindent\textbf{注 3.5.}\ \textit{对几何内窥，我们须以类似方式考察 Hitchin 纤维化的奇异纤维。例如取 $w=e_1$，于是纤维 $\boldsymbol{F}_{e_1}$ 分裂为由比值 $b_2/b_3$ 定义的分支 $\boldsymbol{F}^\pm_{e_1}$，如式 (3.38)。把两个分支在射影空间中紧化，可以看到 $b_4=0$ 的四个点中有两个位于 $\boldsymbol{F}^+_{e_1}$、两个位于 $\boldsymbol{F}^-_{e_1}$。（这其实只用到四个点构成 $Q$ 的轨道、且 $Q$ 交换两分支这一事实。）若限制到 $\tilde{T}^*\mathcal{M}$，两条曲线 $\boldsymbol{F}^\pm_{e_1}$ 在两个相关临界 $z$ 值附近的行为恰如上一段所述。于是每个分数 A-膜都有两个具有这种行为方式的点。}
\par\medskip

\subsubsection{幂零 Higgs 场}

有一种简单的方式来刻画 $\mathcal{M}$ 中这样的点：把 Hitchin 纤维化的纤维 $\boldsymbol{F}$ 视为 $T^*\mathcal{M}$ 中的曲线，这些点正是使该曲线趋于无穷的点。它们对应于存在非零幂零 Higgs 场（其极点留数固定 $E$ 的抛物结构）的稳定抛物丛 $E$。这与刻画几何 Langlands 对偶中出现的扭 D-模有关（见 [La2]）。

例如考虑 $SL_2$ 模空间的 improper 分支。如我们在式 (3.32) 中算出的，Higgs 场 $\varphi=\sum a_iu_i$ 满足 $\Tr\varphi^2=\sum_ia_i^2(x-e_i)$。因此若 $0=\sum_ia_i^2=\sum_ie_ia_i^2$——这正是使 $\boldsymbol{F}$ 中的点不属于 $T^*\mathcal{M}$ 的条件——则 $\Tr\varphi^2$ 恒为零，$\varphi$ 幂零。

对 proper 分支，若 $E=\mathcal{L}\oplus\mathcal{L}^{-1}$，则 Hitchin 纤维 $\boldsymbol{F}$ 在 $\tilde{T}^*\mathcal{M}$ 中趋于无穷的条件是 $\mathcal{L}$ 为 2 阶。（这在式 (3.12) 所用的原始坐标中是清楚的。$\mathcal{L}$ 为 2 阶意味着 $x_0=\infty$ 或 $f(x_0)=0$；对这种 $x_0$，在 $w_0$ 固定时 $h_0\to\infty$。）另一方面，若给 $E$ 赋予使其稳定的抛物结构，则除非 $\mathcal{L}$ 为 2 阶，其极点（位于 $p$）保持该抛物结构的 Higgs 场不可能幂零。事实上，若 $\mathcal{L}$ 不是 2 阶，这种 Higgs 场的一般形式由式 (3.8) 与 (3.9) 给出，它们表明 $\varphi$ 处处幂零只有当 $h_0$ 与 $k_0$ 或 $g_0$ 之一同时为零；而此时在 $p$ 处被 $\varphi$ 的极点部分保持的抛物结构对应于一条位于 $\mathcal{L}$ 或 $\mathcal{L}^{-1}$ 中的复直线 $\Psi\subset E_p$，于是 $E$ 不稳定。反之，若 $\mathcal{L}$ 为 2 阶，我们可以构造一个作为 $\mathcal{L}$ 被 $\mathcal{L}$ 扩张的稳定抛物丛 $E$（直和 $\mathcal{L}\oplus\mathcal{L}$ 对任何抛物结构都不稳定），它容许一个极点部分固定该抛物结构的幂零 Higgs 场 $\varphi$。

下面在亏格 $g$ 的曲线 $C$ 上 $SL_2$ 的语境中，对我们刚才看到的行为作一简要解释。（在以下论证中，只要允许 Higgs 场与二次微分有某些极点，便可纳入分歧。）设 $\varphi_1,\dots,\varphi_{3g-3}$ 为 $H^0(C,K\otimes\mathrm{ad}(E))$ 的一组基。于是对某些复数 $a_i$，一般的 Higgs 场为 $\varphi=\sum_ia_i\varphi_i$。因此 $T^*\mathcal{M}$ 在对应于 $E$ 的点处的纤维 $U_E$ 由 $a_1,\dots,a_{3g-3}$ 参数化。
% ==== 原文 PDF 第 39 页（印刷页 39）====
类似地，设 $w_1,\dots,w_{3g-3}$ 为 $H^0(C,K^2)$ 的一组基。Hitchin 纤维化的纤维 $\boldsymbol{F}$ 由 $\Tr\varphi^2=\sum t_iw_i$（对某些复数 $t_i$）确定。具体地，这给出 $3g-3$ 个二次方程
\begin{equation}\tag{3.50}
P_j(a_1,\dots,a_{3g-3})=t_j,\qquad j=1,\dots,3g-3.
\end{equation}
一般地，这样一组方程恰有 $2^{3g-3}$ 个解（按重数计）。这些就是纤维 $U_E$ 与 $\boldsymbol{F}$ 的交点。若再添加一个变量 $a_0$ 使方程齐次化
\begin{equation}\tag{3.51}
P_j(a_1,\dots,a_{3g-3})=t_ja_0^2,\qquad j=1,\dots,3g-3,
\end{equation}
则 $\mathbb{C}\mathbb{P}^{3g-3}$ 中解的个数（按重数计，且有限时）总是恰为 $2^{3g-3}$。仿射方程 (3.50) 的解数少于该值，当且仅当齐次方程 (3.51) 存在 $a_0=0$ 的解。而齐次方程 $a_0=0$ 的解对应于 $\Tr\varphi^2=0$ 的 Higgs 场 $\varphi$，即幂零 Higgs 场。于是，$U_E$ 与 $\boldsymbol{F}$ 的某些交点趋于无穷，当且仅当 $E$ 容许非零幂零 Higgs 场，与我们例子中所见相同。（该例子涉及分歧情形，其中式 (3.50) 的多项式 $P_j$ 非齐次，故式 (3.51) 的齐次形式略有不同。）
% [Tutor Note: 推导补足。齐次化后交点数恒为 $2^{3g-3}$（Bezout 型计数，按重数计），仿射方程解数少掉的部分正是 $a_0=0$ 的齐次解；而 $a_0=0$ 意味着 $\Tr\varphi^2=\sum t_jw_j$ 的齐次部分为零，即 $\Tr\varphi^2=0$，对 $SL_2$ 即 $\varphi$ 幂零。故「缺失点」与「存在非零幂零 Higgs 场」等价。易错：这里数的是 $U_E$ 与 $\boldsymbol{F}$ 在 $\mathbb{C}\mathbb{P}^{3g-3}$ 中的交，而不是仿射空间中的交。]

\subsubsection{两个分支的比较}

我们所了解到的内容使我们能够找出 $\mathcal{M}_H$ 两分支之间的显式映射，正如 §3.6.2 所预言。

我们想显式地建立如下两者之间的映射：由熟悉方程 $P=0$ 给出的 $\mathcal{M}_H(SL_2)$ 的描述，其中
\begin{equation}\tag{3.52}
P=\rho^2+(2u+w)f(u)-\frac{f'(u)^2}{4}
\end{equation}
及其对 improper 分支 $\mathcal{M}_H(SL_2^*)$ 的类似物
\begin{equation}\tag{3.53}
b_1^2+b_2^2+b_3^2=1
\end{equation}

我们已从 §3.15 知道，所要的变量替换保持 $w$ 不变；对 improper 分支，$w$ 为
\begin{equation}\tag{3.54}
w=e_1b_1^2+e_2b_2^2+e_3b_3^2.
\end{equation}
% [原文疑误: p.39 原文作 "from Section 3.15"，但本节之前并无 §3.15（§3 只到 §3.11）；由上下文（两分支映射、变量替换保持 $w$）疑为 §3.6.2，但无法据图像确证具体节号，故照录原文以便核对。]

此外，所要的变量替换与 $\mathcal{M}_H$ 中某个 Zariski 开集到 $\mathcal{M}$ 的投影交换。理由是：§3.6.2 所用的丛 $E$ 的 $\varphi$-不变 Hecke 修改只依赖于由 $\varphi$ 确定的 $E$ 上的抛物结构，而不依赖于 $\varphi$ 本身。因此它与从 $\mathcal{M}_H^0$ 到 $\mathcal{M}$（稳定抛物丛的模空间）的投影交换。

对 proper 分支，$\mathcal{M}$ 由变量 $u$ 参数化；对 improper 分支，类似地由 $z=(1-b_1)/(b_2-\mathrm{i}b_3)$ 参数化。这些参数化在相差分式线性变换的意义下唯一，因此 $z$ 必映为 $(au+b)/(cu+d)$（$a,b,c,d$ 为常数）。这些常数可由下述事实确定：抛物丛 $E$ 容许幂零 Higgs 场（可能在 $p$ 处有极点，其留数保持抛物结构）的条件也被
% ==== 原文 PDF 第 40 页（印刷页 40）====
两分支之间的映射保持。事实上，该映射同时保持 Hitchin 纤维化与投影 $\mathcal{M}_H^0\to\mathcal{M}$，故它把 Hitchin 纤维上的缺失点映为缺失点。

于是两分支之间的映射把四个点 $u=e_1,e_2,e_3,\infty$ 与式 (3.45) 给出的四个 $z$ 值对应起来。模去 $Q$ 的作用，这就确定了 $z$ 与 $u$ 之间的关系为
% [Tutor Note: 推导补足。$z$ 与 $u$ 都是 $\mathcal{M}\cong\mathbb{C}\mathbb{P}^1$ 上的仿射坐标，故二者只差一个分式线性变换；两分支的映射把 Hitchin 纤维上的缺失点映为缺失点，而四个缺失点 $u=e_1,e_2,e_3,\infty$ 正是 $A(z)$ 的四个零点在 $u$ 平面中的像，四个条件恰好唯一确定该分式线性变换（相差 $Q$ 的作用）。]
\begin{equation}\tag{3.55}
u=\frac{e_2\sqrt{e_3-e_1}+e_3\sqrt{e_2-e_1}+(e_1-\sqrt{e_3-e_1}\sqrt{e_2-e_1})\sqrt{e_2-e_3}z}{\sqrt{e_3-e_1}+\sqrt{e_2-e_1}+\sqrt{e_2-e_3}z}.
\end{equation}

$\mathcal{M}_H(SL_2)$ 是由关于 $\rho$ 的二次方程 (3.52) 给出的 $w-u$ 平面的双重覆盖；类似地，$\mathcal{M}_H(SL_2^*)$ 是由关于 $\tilde v$ 的二次方程 (3.46) 给出的 $w-z$ 平面的双重覆盖。初等计算表明，两个二次方程在变量替换
\begin{equation}\tag{3.56}
\tilde v=\rho\frac{du}{dz}\frac{1}{f(u)}-\frac{1}{2A}\frac{dA}{dz},
\end{equation}
下等价，其中 $A$ 由式 (3.47) 定义。

这给出两个空间之间的映射。最后，直接验证可知 $\mathcal{M}_H(SL_2^*)$ 的复辛形式（由式 (3.42) 给出）映为 $\mathcal{M}_H(SL_2)$ 的复辛形式，后者类似地由
\begin{equation}\tag{3.57}
\Omega=\sigma_0\frac{du\wedge d\rho}{f(u)}=\sigma_0\frac{du\wedge d\rho}{\partial P/\partial w}
\end{equation}
给出。
% [Tutor Note: 易错预警。检验两条复辛形式是否互映，须比较 (3.42) 与 (3.57)：二者都在缺失点处出现极点，$dz\wedge d\tilde v$ 与 $du\wedge d\rho$ 的换算依赖 (3.55)、(3.56) 的 Jacobi 因子。注意 (3.57) 中 $\partial P/\partial w$ 充当了 $f(u)$ 的角色，而由 (3.52) 知 $P$ 关于 $w$ 是线性的；核对时还须保证 (3.42) 中的 $\sigma_0$（同一个乘性常数）在换元下不被吸收，符号约定以 §3.9.1 的 $v=-\sigma_0(b_2-\mathrm{i}b_3)$ 为准。]

由于这一同构，在如 §4 那样比较 A-膜与 D-模时，只需考虑 improper 分支。我们从 $\mathcal{M}_H$ 的几何所能得到的关于 A-膜及其相应 D-模的任何信息，对两分支都相同。

\subsection{轨形的镜像对称}

这里我们简要地把所考虑的一些问题放到镜像对称的更一般语境中。

首先注意关于经典几何的如下简单事实。$\mathcal{M}_H(SO_3)$ 有一个 $A_1$ 奇点，它包含在一条由单个 $\mathbb{P}^1$ 在两点处相交而成的纤维中。（这样的奇异纤维有三条，这里只考虑其中一条。）若把 $A_1$ 奇点爆破，我们消除了（全空间的）奇点，并产生第二个 $\mathbb{P}^1$。爆破产生一个新的椭圆纤维化，其奇异纤维现在含有两个 $\mathbb{P}^1$，每个都在两点处相交。全空间现在光滑，但具有刚才描述的那种奇异纤维。这恰是 $\mathcal{M}_H(SL_2)$ 在其特殊纤维附近的几何。

因此结论是：若通过这一爆破改变物理学家所称的「Kähler 模」，就可把 $SO_3$ 的局部几何转换为 $SL_2$ 的局部几何。

更系统地思考这一局部几何。从超凯勒的观点看，$A_1$ 轨形奇点的形变/消解由三个实参数 $\vec\mu$ 控制。若选定其中一个复结构——自然地选取 Hitchin 纤维化在其下全纯的复结构 $I$——则 $\vec\mu$ 分解为控制 Kähler 类的实参数 $\mu_{\mathbb{R}}$ 与控制复结构的复参数 $\mu_{\mathbb{C}}$。
% ==== 原文 PDF 第 41 页（印刷页 41）====
具有相同 $\mu_{\mathbb{C}}$ 而 $\mu_{\mathbb{R}}$ 不同的簇双有理等价。只有当 $\vec\mu=0$ 时才出现奇点。

在二维 sigma 模型的语境中还有另一个参数，一个类 theta 角；在 [GW] 的 §6.2（该处综述了这一几何）中它被称为 $\eta$，取值于 $\mathbb{R}/\mathbb{Z}$。从其中一个复结构（比如 $I$）的观点看，$\eta$ 与 $\mu_{\mathbb{R}}$ 合并为第二个复参数 $\mu'_{\mathbb{C}}=\mu_{\mathbb{R}}+\mathrm{i}\eta$。该参数控制与该奇点相联系的复化 Kähler 类，而 $\mu_{\mathbb{C}}$ 控制复结构。

sigma 模型奇异当且仅当 $\vec\mu=\eta=0$，等价地 $\mu_{\mathbb{C}}=\mu'_{\mathbb{C}}=0$。sigma 模型具有奇点是内在条件，因而在任何对偶下不变。在 Type IIA 超弦理论的语境中，sigma 模型在 $\vec\mu=\eta=0$ 处的奇点导致时空中的 $SU(2)$（或 $A_1$）规范对称；这或许是关于群 $A_1$ 与 $A_1$ 奇点之间关系的最令人满意的陈述。

在几何 Langlands 中，我们希望把 $\mathcal{M}_H(SO_3)$ 的 sigma 模型中的 $\mathbb{R}^4/\mathbb{Z}_2$ 奇点当作轨形点处理。按 [As]，恰在 $\vec\mu=0$、$\eta=1/2$ 时 sigma 模型可当作轨形处理。在此值处 $\mu_{\mathbb{C}}=0$、$\mu'_{\mathbb{C}}\ne0$。这意味着离开点 $\mu_{\mathbb{C}}=\mu'_{\mathbb{C}}=0$ 的形变最好理解为消解而非形变，尽管这是非经典意义下的消解，涉及的是 $\eta=\operatorname{Im}\mu'_{\mathbb{C}}$ 而非 $\mu_{\mathbb{R}}=\operatorname{Re}\mu'_{\mathbb{C}}$。相比之下，要得到 $\mathcal{M}_H(SL_2)$ 的局部几何，如开头所解释的，必须作经典消解，即 $\operatorname{Re}\mu'_{\mathbb{C}}\ne0$。
% [Tutor Note: 易错预警。$\vec\mu$ 的三个实分量是 $A_1$ 奇点超凯勒商的三重矩映射值：$\mu_{\mathbb{R}}$ 是 Kähler 类（实），$\mu_{\mathbb{C}}$ 是复结构（复），而 $\eta$ 是二维 sigma 模型的 theta 角，与 $\mu_{\mathbb{R}}$ 组成 $\mu'_{\mathbb{C}}$。因此「消解」与「形变」的区分取决于动的是 $\operatorname{Re}\mu'_{\mathbb{C}}$ 还是 $\operatorname{Im}\mu'_{\mathbb{C}}$；在轨形点 $\eta=1/2$ 处只能作非经典（$\operatorname{Im}$ 方向）消解。切勿把这几个参数与 Hitchin 底 $\boldsymbol{B}$ 的坐标混淆。]

于是在此情形，Hitchin 纤维化各纤维上的 T-对偶把涉及 $\operatorname{Im}\mu'_{\mathbb{C}}$ 的非经典消解与涉及 $\operatorname{Re}\mu'_{\mathbb{C}}$ 的经典消解互换。这与下述事实相符：由于 Hitchin 纤维化在复结构 $I$ 下全纯，该纤维化各纤维上的 T-对偶把复参数（如 $\mu_{\mathbb{C}}$）映为复参数，把 Kähler 参数（如 $\mu'_{\mathbb{C}}$）映为 Kähler 参数。

最后一点评论：在 §5.2 中我们会遇到涉及奇点 $\mathbb{C}^{2n}/\mathbb{Z}_2$（$n>1$）的这一切的类比。此时在经典几何中该奇点既不能形变也不能消解，而在 sigma 模型中它只能作为轨形处理。因此与上述讨论没有密切的类比。

\subsection{与 Seiberg--Witten 理论的关系}

我们这里考虑的 Hitchin 纤维化与 Seiberg--Witten 理论中出现的椭圆纤维化非常相似。例如，模去一些初等的变量替换，$SO_3$ 模空间的方程 (3.39) 与 [SW1] 的式 (16.24) 一致。后一方程描述规范群为 $SU(2)$ 的四维 $\mathcal{N}=2^*$ 理论\footnote{这是带伴随表示中一个大规模超多重态的 $\mathcal{N}=2$ 理论。当超多重态质量变大时，该理论可约化为极小的 $\mathcal{N}=2$ 理论，后者是 Seiberg--Witten 理论的基本情形。}的 Seiberg--Witten 纤维化。我们在此与内窥相联系的 Hitchin 纤维化的同一批奇异纤维，在 Seiberg--Witten 理论中也起核心作用。物理上，它们描述无质量带电超多重态的出现；数学上，它们给出 Seiberg--Witten 理论应用于四维流形时的主要贡献。
% [背景补充: $\mathcal{N}=2^*$ 是 $\mathcal{N}=2$ 超 Yang--Mills 加上伴随表示中的一个大规模超多重态；当质量趋于无穷时回到纯 $\mathcal{N}=2$ 理论，即 Seiberg--Witten 理论的基本情形。其 Seiberg--Witten 曲线正是 Coulomb 分支上的椭圆纤维化，质量参数对应本文的分歧参数 $\sigma$。]

% ==== 原文 PDF 第 42 页（印刷页 42）====

$\mathcal{N}=2^*$ 理论的 Seiberg--Witten 纤维化最初并未被理解为 Hitchin 纤维化，但后来 [DW] 做到了这一点，作为理解 $SU(N)$ 的 $\mathcal{N}=2^*$ 理论的一步。

同一椭圆纤维化出现在两个不同问题中，这一事实可以通过把一个特定 M-理论五膜（fivebrane）的六维世界体积以两种不同方式从六维约化到四维来解释。考虑 11 维流形 $X=\mathbb{R}^2\times\mathbb{T}^2\times T^*C\times\mathbb{R}^3$ 上的 M-理论，其中 $C$ 是椭圆曲线。令 $D\subset T^*C$ 为与 $C$ 上带一个分歧点的 $SL_2$ Higgs 丛（对某个 $w$ 值）相联系的谱曲线。\footnote{这一概念在 §4.1 中综述。} 并考虑世界体积为 $\mathbb{R}^2\times\mathbb{T}^2\times D$ 的 M-理论五膜。有两种方式看待这一情形，取决于 $\mathbb{T}^2$ 还是 $C$ 更小：
\begin{enumerate}[label=(\arabic*),leftmargin=2.2em,itemsep=2pt,topsep=2pt]
\item 若 $\mathbb{T}^2$ 较小，第一步是考虑五膜在 $\mathbb{T}^2$ 上的紧化。这给出规范群为 $U(2)$ 的 $\mathcal{N}=4$ 超 Yang--Mills 理论（因为投影 $D\to C$ 一般是一个双重覆盖，规范群为 $U(2)$，或有效地为 $SU(2)$，中心作用甚微）。几何 Langlands 与随后在 $C$ 上紧化到二维相联系；特定的曲线 $D$ 描述了在与分歧相联系的面算子存在时的行为，如 [GW] 的 §6.4 所述。
\item 若 $C$ 较小，先考虑在 $C$ 上的紧化，[Wi3] 指出这给出 $\mathbb{R}^2\times\mathbb{T}^2$ 上的 $\mathcal{N}=2^*$ 理论。椭圆纤维化 (3.39) 描述了 $\mathbb{R}^4$ 上的该理论。进一步从 $\mathbb{R}^4$ 紧化到 $\mathbb{R}^2\times\mathbb{T}^2$ 的结果是，该椭圆纤维化必须被赋予一个超凯勒度规 [SW2]。
\end{enumerate}
% [Tutor Note: 物理图像。两个紧化给出同一椭圆纤维化：先缩 $\mathbb{T}^2$ 得到 $U(2)$ 的 $\mathcal{N}=4$ 超 Yang--Mills（其 Hitchin 纤维化即 $D\to C$），先缩 $C$ 得到 $\mathcal{N}=2^*$ 的 Seiberg--Witten 纤维化；二者在 $\mathbb{R}^2\times\mathbb{T}^2$ 上相会，说明 Hitchin 纤维化与 SW 椭圆纤维化是同一个六维对象的两个面。$\mathcal{N}=2^*$ 中的超多重态质量对应本文的分歧参数 $\sigma$（或 $w$）。]
